\section{Data Preparation}\label{data-preparation}

This chapter covers the CRISP-DM data preparation phase for the frozen analytics base table (ABT). Chapter 4 described the condition of the data. This chapter explains what was done to turn that file into a modeling-ready dataset. The focus here is on filtering, missing-data treatment, feature engineering, encoding, and the final train-test split used by the modeling script.

\subsection{\texorpdfstring{\textbf{Scope of Preparation Work}}{Scope of Preparation Work}}\label{scope-of-preparation-work}

The preparation work started from the frozen ABT with 2{,}047 rows and 50 columns. The ABT already reflected the six-LGU study scope used in this thesis: Cebu City, Mandaue City, Lapu-Lapu City, Talisay City, Minglanilla, and Consolacion. Naga City stayed in the system only as a CBD distance anchor through dist\_naga\_city\_m. Its four scraped listings were excluded before the ABT freeze because they were too few to support stable LGU-level learning.

The goal of this phase was more specific than general ABT cleaning. It was to build the exact dataset used by the modeling pipeline in thesis\_main/Scripts/run\_models.py. This meant keeping the ABT as the full reference file, but applying a narrower set of preparation steps for the training data.

The preparation flow is summarized in Table~\ref{tab:chapter5-prep-flow}. The biggest change was the restriction to the open\_market segment. After that, the script removed only rows that could not support the target variable or the spatial-lag feature.

\begin{table}[htbp]
\caption{Preparation Flow from Frozen ABT to Modeling-Ready Dataset}
\label{tab:chapter5-prep-flow}
\begin{center}
\begin{tabular}{p{0.34\textwidth} p{0.12\textwidth} p{0.12\textwidth} p{0.17\textwidth}}
\toprule
Preparation Step & Rows Kept & Share of ABT & Rows Removed \\
\midrule
Frozen ABT before Chapter 5 filtering & 2{,}047 & 100.00\% & -- \\
Filter to open\_market & 1{,}619 & 79.09\% & 428 \\
Drop three integrity rows & 1{,}616 & 78.95\% & 3 \\
Drop null price\_per\_sqm & 1{,}509 & 73.72\% & 107 \\
Drop null spatial\_lag\_price & 1{,}491 & 72.84\% & 18 \\
\bottomrule
\end{tabular}
\end{center}
\end{table}

By the end of this phase, 1{,}491 rows remained for model fitting. This means 72.84\% of the frozen ABT entered the final training workflow. The removed rows were not dropped because they were statistically inconvenient. They were dropped because they could not support the definition of the target variable or the required neighborhood context.

\begin{center}\rule{0.5\linewidth}{0.5pt}\end{center}

\subsection{\texorpdfstring{\textbf{Cleaning and Filtering}}{Cleaning and Filtering}}\label{cleaning-and-filtering}

The first cleaning step was a scope decision. The ABT kept three market segments for reference: open\_market, bank\_ropa, and floor\_price. However, the deployed task of the thesis was the estimation of open-market residential prices. Because of that, the modeling script filtered the ABT to market\_segment = open\_market before any encoding or model fitting. This removed 428 rows from the broader file and left 1{,}619 open-market observations.

This step also aligned the training data with the thesis valuation basis. Open-market listings represent the closest available proxy to Market Value in the IVS sense, while bank ROPA and administrative floor-price records reflect different pricing conditions and should not be pooled into the same training target \parencite{ivs2025}. The ABT still retained those other segments, but they were treated as reference layers rather than training labels.

A second cleaning step handled obvious integrity failures. Three rows with property\_id values 468, 714, and 769 were removed before model fitting. These rows had implausible unit prices that were not defensible for residential valuation in Metro Cebu. Their removal was a hard data-integrity action, not a general outlier-pruning rule.

The remaining row filters were tied to model requirements. Rows with null price\_per\_sqm were removed because the target variable could not be computed. Rows with null spatial\_lag\_price were also removed because the neighborhood-price feature was part of the final design. After these two filters, the open-market subset decreased from 1{,}616 to 1{,}491 rows.

Property-type cleanup was already applied before this stage. The ABT used a standardized six-class residential taxonomy: Condominium, House and Lot, Townhouse, Single Detached, Vacant Lot, and Apartment. The earlier generic Residential label had already been removed before this stage, so Chapter 5 started from a residential-only file with the cleaned final taxonomy.

\begin{center}\rule{0.5\linewidth}{0.5pt}\end{center}

\subsection{\texorpdfstring{\textbf{Missing Data Treatment}}{Missing Data Treatment}}\label{missing-data-treatment}

Missing-data treatment followed the practical pattern of the project: keep the file usable, preserve meaning when possible, and avoid inventing information that the sources did not actually provide.

The first structural decision involved size variables. The thesis used area\_sqm as the main size feature, with floor\_area\_sqm preferred when available and lot\_area\_sqm used as fallback. This was paired with an is\_vacant\_lot flag so that land-only listings stayed visible to the model. This design was important because vacant lots are priced on land area, while most open-market condominium listings do not carry a meaningful lot-area field.

The second issue was structural absence in lot\_area\_sqm. In the open-market segment, lot\_area\_sqm was fully missing because Lamudi listings did not systematically provide lot area. The modeling script treated this as an all-null column, filled it with 0 inside the feature matrix, and then dropped log\_lot\_area\_sqm from the OLS matrix because it became a constant field. This kept the matrix complete without pretending that lot area was observed for those listings.

The ABT already carried preparation flags for bedrooms, bathrooms, and floor area. In the open-market segment, the file contained 449 rows with bedrooms\_imputed = 1, 299 rows with bathrooms\_imputed = 1, and 107 rows with floor\_area\_imputed = 1. After the final training filters, the remaining 1{,}491-row modeling subset still contained 397 bedroom-imputed rows and 253 bathroom-imputed rows, while floor\_area\_imputed dropped to 0 because the rows with unresolved floor-area issues were removed together with the null target rows.

The final model-build script still applied one more safeguard step. After filtering and encoding, it checked the remaining numeric columns for nulls and used median imputation where needed. In the current frozen pipeline, this affected bedrooms, bathrooms, bir\_zonal\_cr\_median, and bir\_zonal\_rc\_median. The all-null lot\_area\_sqm column was handled separately and filled with 0 rather than sent through the median imputer.

Table~\ref{tab:chapter5-missing-treatment} summarizes the main missing-data actions that carried into the final matrix.

\begin{table}[htbp]
\caption{Main Missing-Data Treatments Used in Chapter 5}
\label{tab:chapter5-missing-treatment}
\begin{center}
\small
\begin{tabular}{p{0.19\textwidth} p{0.23\textwidth} p{0.19\textwidth} p{0.18\textwidth}}
\toprule
Field or Issue & Treatment Used & Why It Was Needed & Effect on Final Matrix \\
\midrule
area\_sqm / is\_vacant\_lot & Use floor area first, then lot area fallback; keep vacancy flag & Separate built space from land-only listings & Kept one usable size signal across property types \\
lot\_area\_sqm in open\_market & Fill all-null column with 0 in the model script & Lamudi rows structurally lacked lot-area values & Matrix stayed complete; constant log form was dropped from OLS \\
bedrooms and bathrooms & Keep audit flags in the ABT, then median-impute remaining numeric nulls & Interior details were incomplete in part of the listings & Structural fields stayed usable without dropping many rows \\
bir\_zonal\_cr\_median and bir\_zonal\_rc\_median & Median-impute remaining numeric nulls after filtering & Some zonal classes were still missing by location & Administrative columns stayed usable in the final matrix \\
\bottomrule
\end{tabular}
\end{center}
\end{table}

No separate feature scaling or normalization step was applied in this phase. The pipeline did not include a scaler, and the final models did not require one for the implemented workflow.

\begin{center}\rule{0.5\linewidth}{0.5pt}\end{center}

\subsection{\texorpdfstring{\textbf{Market Segment Treatment}}{Market Segment Treatment}}\label{market-segment-treatment}

The market\_segment field stayed in the ABT because it was useful for audit and provenance work. It allowed the project to keep open\_market, bank\_ropa, and floor\_price rows in one canonical file. However, it was not used as a model feature in the final training workflow.

Once the dataset was filtered to open\_market only, market\_segment became a constant column. A constant string field adds no useful variation, so it was excluded from the feature matrix. This prevented the model from learning a meaningless dummy for a segment that no longer varied in the training data.

This step also kept the model aligned with the actual prediction objective of the Streamlit app. The deployed map estimates open-market residential price levels across Metro Cebu. It does not estimate forced-sale or administrative floor-price values.

\begin{center}\rule{0.5\linewidth}{0.5pt}\end{center}

\subsection{\texorpdfstring{\textbf{Outlier Flagging}}{Outlier Flagging}}\label{outlier-flagging}

The project treated outliers in two different ways. First, the ABT kept a soft price\_outlier\_flag field. This flag marked values outside the 1st to 99th percentile range of price\_php in the full ABT. In the frozen file, 19 rows had price\_outlier\_flag = True: 14 in open\_market, four in bank\_ropa, and one in floor\_price. These rows were not automatically removed just because they were extreme.

For bank ROPA rows, the same percentile rule was backfilled in the ABT so the flag would be consistent across sources. This preserved the flag as a diagnostic variable without forcing a blanket deletion policy.

Second, the project used a much narrower hard-drop rule for obvious integrity failures. Only the three implausible open-market rows with property\_id 468, 714, and 769 were removed before model fitting. This distinction mattered. Soft outlier flags were used to warn the analyst that a price was unusual. Hard integrity drops were reserved for rows that were not credible enough to keep in a residential valuation model.

\begin{center}\rule{0.5\linewidth}{0.5pt}\end{center}

\subsection{\texorpdfstring{\textbf{Feature Engineering}}{Feature Engineering}}\label{feature-engineering-ch5}

Most of the heavy feature engineering had already been completed before this chapter, but Chapter 5 still had to decide which engineered fields would actually enter the model matrices.

The project kept the eight network-distance variables for the tree-based models: Cebu Business Park, Mandaue CBD, Mactan CBD, SRP, Talisay Tabunok, Consolacion, Naga City, and Airport distances. The is\_mactan\_island flag was also retained in the full matrix because Lapu-Lapu properties sit in a different road-network and bridge-access context.

For accessibility, the final ABT retained the nine category-level Metro Cebu Residential Accessibility Index (MCRAI) fields. The older amenity\_score\_* fields had already been dropped because they duplicated the same concept with a weaker formulation. The ABT also stored mcrai\_composite, but the model matrix excluded it because it is a weighted summary of the nine component scores and would create redundancy if used together with them.

The spatial\_lag\_price field was treated as a required engineered feature, not an optional add-on. Because of that, rows with null spatial\_lag\_price were removed during the filtering stage instead of being imputed.

The OLS baseline used a narrower version of the feature set than Random Forest and XGBoost. This was done to reduce collinearity and coefficient instability. In the current script, X\_ols dropped dist\_talisay\_tabunok\_m, is\_mactan\_island, bir\_zonal\_rr\_log, is\_vacant\_lot, latitude, longitude, lot\_area\_sqm, floor\_area\_sqm, area\_sqm, and log\_lot\_area\_sqm. It then added log\_floor\_area\_sqm as the area term for the log-linear hedonic baseline.

\begin{center}\rule{0.5\linewidth}{0.5pt}\end{center}

\subsection{\texorpdfstring{\textbf{Encoding and Final Modeling-Ready Dataset}}{Encoding and Final Modeling-Ready Dataset}}\label{encoding-and-final-modeling-ready-dataset}

The final preparation step converted the cleaned dataset into model matrices. The script one-hot encoded city and property\_type using drop\_first = True. This produced five city dummy variables and five property-type dummy variables in the full feature matrix.

After encoding and numeric-null handling, the full matrix used for Random Forest and XGBoost had shape 1{,}491 $\times$ 48. The OLS baseline matrix had shape 1{,}491 $\times$ 39 after the extra stability drops described earlier. Table~\ref{tab:chapter5-final-matrix} summarizes the result.

\begin{table}[htbp]
\caption{Final Modeling-Ready Matrices and Split}
\label{tab:chapter5-final-matrix}
\begin{center}
\begin{tabular}{p{0.28\textwidth} p{0.18\textwidth} p{0.34\textwidth}}
\toprule
Output & Shape or Size & Notes \\
\midrule
Full tree-model matrix (X\_full) & 1{,}491 $\times$ 48 & Used for Random Forest and XGBoost \\
OLS matrix (X\_ols) & 1{,}491 $\times$ 39 & Used for the log-linear baseline after extra stability drops \\
Training split & 1{,}192 rows & 80\% of the modeling-ready data \\
Test split & 299 rows & 20\% of the modeling-ready data \\
Categorical encoding & 5 city dummies, 5 property-type dummies & Created with one-hot encoding and drop\_first = True \\
\bottomrule
\end{tabular}
\end{center}
\end{table}

The final modeling-ready subset remained concentrated in the same main residential markets seen in Chapter 4, especially Lapu-Lapu City and Cebu City. Even after the row drops, the dataset still covered all six LGUs and all six standardized residential property types.

Overall, the data preparation phase was conservative. It did not chase a perfectly tidy dataset. Instead, it applied a small number of defensible filters, preserved the main valuation signals, and made the feature matrix complete enough for model fitting. Chapter 6 discusses how these prepared matrices were used in OLS, Random Forest, and XGBoost.
